\chapter*{ABSTRACT}
\addcontentsline{toc}{chapter}{Abstract}

Software organisations run each project through separate tools for version
control, issue tracking, static code analysis and CI/CD, and each tool reports
its own metrics on its own scale. Judging whether a project is healthy therefore
means reconciling four or five dashboards by hand. At Spectrum Software \&
Consulting Pvt.\ Ltd.\ (SSCL), the industry partner that raised the problem,
this is done rarely. Projects cannot be compared like for like, and a slowly
deteriorating project is noticed only once deadlines slip. The tools also
capture nothing about how the team is faring, or about what management did in
response and whether it helped.

This project designed and built \emph{Pulse}, a multi-tenant engineering-health
platform. Scheduled and on-demand synchronisation pulls metrics through a
connector abstraction into immutable snapshots. A deterministic rule-based model
converts them into seven health scores and an overall score on a common 0--100
scale. An absent metric is excluded and the remaining weights are renormalised,
and every score decomposes into the signals, weights and contributions that
produced it. AI-generated surveys, delivered by email, collect team feedback
without storing any respondent identity. An action log records management
interventions and prompts a later review of their effectiveness. The system is
built with TypeScript, PostgreSQL, Redis and React. It is deployed on free-tier
services at no monetary cost.

Verification combined 142 automated tests enforced as a required CI check,
inspection, end-to-end demonstration on the team's own repositories, and four
targeted investigations. The system was also reviewed with SSCL's Chief
Technical Officer in five meetings. Of 41 functional requirements, 38 are met;
of 26 non-functional requirements, 21 are met. The investigations showed that
renormalisation stops a project being penalised for an unconnected tool (a
score of 80 rather than 12 in a worked example). They also found that no stored
record links a survey response to its author. Six connectors across the four
tool categories were delivered, two more than planned. The main limitations are
that AI-assisted action search does not fall back automatically when its
provider is unavailable, and that the rate-limit capacity per shared credential
is unquantified. Overall test coverage is also low, with no frontend tests.

\vspace{12pt}
\noindent\textbf{Keywords:} software project health; engineering metrics;
explainable scoring; anonymous surveys; tool integration
\clearpage
